\section{Discussion}

\paragraph{Guideline.}
The main takeaway is not that DFL is broadly unnecessary, but that its room for
improvement depends on how prediction errors affect the downstream objective.
DFL is most useful when prediction errors change the objective value of the
selected solution, not merely its identity. If a predictor selects a different
feasible solution whose true value is nearly oracle-level, then the normalized
regret left for any decision-aware training method to remove is small.

\paragraph{When two-stage may be enough.}
In problems with close substitute solutions, a two-stage model can have decision
performance similar to DFL even when it does not recover the oracle solution
identity. This is especially plausible in decomposable packing-style feasible
sets, where a solution is assembled from multiple components and one component
can often be exchanged for another with little loss in objective value. Our KEP
rank-2 evidence, cycle-length sensitivity, density perturbations, and case
studies all point to this mechanism: many alternative packings remain close to
the oracle under the true objective.

\paragraph{When DFL may help.}
The contrast is that some feasible sets amplify local prediction errors. In
coupled or path-like problems, an incorrect local preference can redirect the
whole decision rather than swapping one nearly equivalent component. In that
setting, selected-solution identity and objective value are more tightly linked,
and DFL can have substantially more room to improve over a two-stage baseline.
This helps explain why shortest-path benchmarks can show large SPO+ gains while
KEP-like packing problems may show more similar performance between two-stage
and optimization-aware training.

\paragraph{Limitations.}
Our KEP experiments cover one synthetic regime, so the numerical magnitudes
should not be read as universal constants for kidney exchange. The density
sensitivity analysis uses selected case graphs rather than a full distributional
benchmark, and the randomized toy experiments are controlled mechanism studies
rather than full benchmark evidence. They also clarify that the claim is
conditional: packing and stable-set problems are not automatically easy for
two-stage models. When prediction noise selects many wrong local components, or
when alternatives are well separated, regret can increase and DFL has more room
to help. The evidence supports the proposed explanation, but broader benchmarks
are needed before treating Property X as a complete taxonomy of when DFL helps.

\paragraph{Future work.}
A natural next step is to test Property X directly on standard stable-set, set
packing, matching, and knapsack-style benchmarks. The key empirical question is
whether best/second-best diagnostics can predict when two-stage and DFL methods
will have similar decision performance. If this diagnostic transfers across
problem classes, it could become a practical pre-check for whether the extra
complexity of decision-focused training is likely to pay off.
